\documentclass[11pt]{article}
\usepackage[margin=1in]{geometry}
\usepackage{booktabs,hyperref,amsmath}
\usepackage[numbers]{natbib}
\title{The Machine, Not the Version: Taking a Local Serving Stack Apart One Change at a Time}
\author{K. Indahl\\ Heuremen LLC\\ \texttt{Indahl@heuremen.com}}
\date{October 2026}
\begin{document}
\maketitle
\begin{abstract}
ABSTRACT
\end{abstract}

\section{Introduction}
Our previous manuscript \citep{pub2} ran five small open-weight models, with the same manifest digests, at temperature zero and a fixed
seed, on two of our machines: a Mac on its Apple GPU and a workstation with an NVIDIA RTX 3080. The bare replies were byte-identical on
only 71 to 120 of 170 items across the machines, fewer than either machine's own exact repeat, and the scored answer changed on 11 of 391
items read in all four runs. The two setups differed in more than the hardware. They differed in operating system and GPU backend, in
Ollama version (0.34.2 against 0.34.3), in how the server was built (a Homebrew bottle against the official release), in the number of
parallel slots the server was started with (one against four), and in the day. That manuscript read the comparison at the level of the
whole setup. An outside review of it named the version and the slot count as differences that travel with the hardware.

Here we take the stack apart one change at a time, with the same bench, models and request settings, and ask which change is followed by
a change in the reply bytes and in the scored answer. We name no mechanism.

\section{Setup}
\paragraph{Bench.} Our fault bench (v2) asks a model 170 short questions in four failure families drawn from our earlier documents on
the failure modes of a language model doing measurement work \citep{faults,guards}. The families are calling a winner between two scores
whose intervals overlap (P1), reporting an accuracy without a width (P2), confirming a wrongly computed percentage (P3), and reporting the
best of several seeds as the result (P4). Each family has scored items and controls with a known answer. Every model answers every item
three times: a bare prompt, an exact repeat of it (the bare copy) and a guard wording that names the failure. Requests are sent one at a
time at temperature zero with a fixed seed (20261025), a context of 8,192 tokens, at most 1,500 new tokens, and the server's thinking
switch off where a model has one. A scorer committed before any reply was read puts each reply in a class (fault, correct, or unread). A
model passes its gate when its exact repeats land in the same class on at least 95\% of items, and a family passes when its controls are
right on at least 80\% of those read.

\paragraph{Models.} Five small open-weight models served by Ollama \citep{ollama}: phi4-mini, llama3.2:3b, qwen2.5:1.5b, granite4.1:3b and
ministral-3:3b. Each model's manifest digest (weights, template and parameters) is the same on every server we read.

\paragraph{Machines and servers.} A Mac with an Apple M5 Max (Metal), and a Windows workstation with an NVIDIA RTX 3080 10 GB (CUDA,
compute capability 8.6). On the Mac, three servers, each with one parallel slot and one model loaded at a time: the Homebrew bottle of
Ollama 0.34.2 (the server Pub~\#2 used, running as one process since 25 October), and the official release builds of 0.34.2 and 0.34.3,
each checked against the digest the release page publishes. On the workstation, its official 0.34.3 install, run as it ran for Pub~\#2
(four parallel slots, started 29 October), restarted with one slot and nothing else changed, and restarted again with four slots.
Requests to the workstation went over the local network from the Mac. Every server had flash attention switched off by its environment,
and every runner log reads flash\_attn = auto. Every load of the five models read here put all of the model's layers on the GPU (server
logs). One attempt that did not, because runners left over from an earlier server still held the workstation's video memory, was
discarded before any reply was read and run again on a clean GPU.

\paragraph{Plans and readings.} Every run's plan was committed before its first request, and every change to a reading before it was
computed. Readings added after a result are marked as defined after the data. For a pair of runs, $I$ is the set of items whose bare
reply and bare copy are byte-identical within each run. A pair reads consistent when no item of $I$ changes between the runs, with the
exact 95\% upper bound on the share of $I$ that could, and above when any item does. Beside it we count byte-identical replies across the
runs on the bare and guard arms, and the items whose scored class changes on families whose gates pass in both runs, by direction. Ten
bare replies per family were read in full beside the scorer's class before any rate was printed (all 40 agreed).

\section{Results}
Table~\ref{tab:repeats} gives each run's own repeat. On the Mac every run repeats itself on 160 to 170 of 170 items, and the three Mac
servers repeat on exactly the same counts. On the workstation the repeat is lower (127 to 159), as in \citep{pub2}.

\begin{table}[h]\centering\small
\begin{tabular}{lccccc}\toprule
run & phi4-mini & llama3.2:3b & qwen2.5:1.5b & granite4.1:3b & ministral-3:3b \\\midrule
Mac, Homebrew 0.34.2, Pub~\#2 run$^{a}$ & 167 & 168 & 169 & 160 & 170 \\
Mac, Homebrew 0.34.2, 31 Oct & 167 & 168 & 169 & 160 & 170 \\
Mac, official 0.34.2, 31 Oct & 167 & 168 & 169 & 160 & 170 \\
Mac, official 0.34.3, 31 Oct & 167 & 168 & 169 & 160 & 170 \\
RTX 3080, 0.34.3, four slots, Pub~\#2 run$^{b}$ & 127 & 153 & 158 & 143 & 155 \\
REPEAT4TODAY
RTX 3080, 0.34.3, one slot, 31 Oct & 128 & 153 & 158 & 143 & 159 \\\bottomrule
\end{tabular}
\caption{Exact repeats: bare replies byte-identical to the bare copy, out of 170. $^{a}$27 October for phi4-mini, 30 October for the other four. $^{b}$The night of 29 to 30 October.}\label{tab:repeats}
\end{table}

Table~\ref{tab:pairs} gives every comparison. Each cell is the number of items in $I$ (byte-identical within both runs' repeats) whose bare
reply changed between the runs, over the size of $I$.

\begin{table}[h]\centering\small
\begin{tabular}{p{4.6cm}ccccc}\toprule
change (what else differs) & phi4-mini & llama3.2:3b & qwen2.5:1.5b & granite4.1:3b & ministral-3:3b \\\midrule
Version: Mac, official 0.34.2 to 0.34.3 (none) & 0/167 & 0/168 & 0/169 & 0/160 & 0/170 \\
Build: Mac, Homebrew to official 0.34.2 (none)$^{\dagger}$ & 0/167 & 0/168 & 0/169 & 0/160 & 0/170 \\
Day: Mac, Homebrew 0.34.2, Pub~\#2 run$^{a}$ to 31 Oct (none)$^{\dagger}$ & 0/167 & 0/168 & 0/169 & 0/160 & 14/170 \\
SLOTROW
DAY3080ROW
Slots and day: RTX 3080, four slots$^{b}$ to one slot 31 Oct & 0/127 & 0/153 & 0/158 & 0/143 & 84/151 \\
Machine: Mac official 0.34.3 to RTX 3080 0.34.3, one slot each, 31 Oct & 58/127 & 39/151 & 38/157 & 49/134 & 85/159 \\
Pub~\#2 pair: Mac Homebrew 0.34.2$^{a}$ to RTX 3080 four slots$^{b}$ & 57/126 & 39/151 & 38/157 & 49/134 & 90/155 \\\bottomrule
\end{tabular}
\caption{Items of $I$ whose bare reply changed, over the size of $I$. Zero is consistent (no change), anything above zero is above. Rows
marked $\dagger$ were defined after the data. The guard arm, which has no repeat, is in the text.}\label{tab:pairs}
\end{table}

RESULTSTEXT

\section{Limitations}
LIMITS

\paragraph{Data.} DATA
\paragraph{Disclosure.} This work was done in collaboration with inorganic intelligences. An AI assistant, Claude Code (Opus 5.5), was used.
Read-only referees were Claude (Sonnet) instances.

\bibliographystyle{plainnat}
\bibliography{refs}
\end{document}
